\documentclass[11pt]{article}
\usepackage[T1]{fontenc}
\usepackage{lmodern}
\usepackage[margin=1in]{geometry}
\usepackage{amsmath,amssymb,amsthm}
\usepackage{booktabs,array}
\usepackage[hidelinks]{hyperref}

\hypersetup{
  pdftitle={Left-Right Filling Permutations and the 1-Positions of the Morphism 0 -> 001, 1 -> 101: Proof and Lean 4 Formalization of Kimberling's OEIS A189640/A026232 Conjecture},
  pdfauthor={Derek Blair}
}

\newtheorem{theorem}{Theorem}[section]
\newtheorem{proposition}[theorem]{Proposition}
\newtheorem{lemma}[theorem]{Lemma}
\newtheorem{corollary}[theorem]{Corollary}
\theoremstyle{definition}
\newtheorem{definition}[theorem]{Definition}
\newtheorem{remark}[theorem]{Remark}

\newcommand{\N}{\mathbb{N}}
\newcommand{\Z}{\mathbb{Z}}
\newcommand{\Q}{\mathbb{Q}}
\newcommand{\seq}[1]{\texttt{#1}}
\newcommand{\OnePos}{\mathrm{OnePos}}
\newcommand{\ZeroPos}{\mathrm{ZeroPos}}

\setlength{\emergencystretch}{2em}

\title{Left-Right Filling Permutations and the $1$-Positions of $0 \to 001,\, 1 \to 101$:\\[2pt]
\large Proof and Lean~4 Formalization of Kimberling's OEIS \seq{A189640}/\seq{A026232} Conjecture}
\author{Derek Blair}
\date{24 November 2026}

\begin{document}
\maketitle

\begin{abstract}
Let $a = \seq{A189640}$ be the $1$-based infinite binary fixed point of the $3$-uniform morphism $\alpha(0) = 001$, $\alpha(1) = 101$ starting from seed $0$, and let $\Pi_{66} = \seq{A026166}$ and $\Pi_{42} = \seq{A026142}$ be Clark Kimberling's permutations of the positive integers generated by the left-right filling algorithms with step offsets $h_{66}(n) = \lfloor (n-1)/2 \rfloor$ and $h_{42}(n) = \lfloor (n+1)/2 \rfloor$, respectively. Let $s = \seq{A026231}$ enumerate the set $\{K \ge 1 : \Pi_{66}(K) = \Pi_{42}(K) + 1\}$ in increasing order, and let $\seq{A026232}(n) = \frac{1}{3}(s(n) + 1)$. On April~24, 2011, Kimberling asked in OEIS~\seq{A189640} whether the $0$-positions of $\seq{A189640}$ are given by \seq{A026138} and whether its $1$-positions are given by \seq{A026232}. In 2020, F.~M.\ Dekking (\emph{J.\ Integer Seq.} \textbf{23}, Article~20.4.7) settled the companion $0$-positions question and characterized \seq{A026142}, while explicitly leaving the permutation \seq{A026166} ``to the interested reader.''

We close that gap and resolve Kimberling's $1$-positions conjecture unconditionally for all $n \ge 1$. Under a uniform regularity criterion on step offsets, we establish no-overwrite persistence for left-right filling algorithms and derive a complete residue-class characterization of \seq{A026166}: for all $r \ge 1$, $\Pi_{66}(3r) = 6r - 2$, $\Pi_{66}(3r + 1) = 2r + 1$, and $\Pi_{66}(3r + 2)$ equals $6r + 2$ if $\Pi_{66}(r + 2) = 2r + 2$ and $2r + 2$ otherwise. By strong induction on $t \in \N$, the branching predicate $\Pi_{66}(t + 2) \ne 2t + 2$ satisfies the exact same $3$-way ternary recurrence as $\seq{A189640}(t + 1) = 1$. Comparing $\Pi_{66}$ against $\Pi_{42}$ modulo~$3$ then yields the exact set equivalence
\[
\Pi_{66}(K) = \Pi_{42}(K) + 1 \iff \exists\, m \ge 1,\; \seq{A189640}(m) = 1 \;\wedge\; K = 3m - 1
\]
for every $K \ge 1$. Commuting order-preserving enumeration along $m \mapsto 3m - 1$ proves both exact division-by-$3$ integrality $\seq{A026231}(n) \equiv 2 \pmod{3}$ and the positional identity $\OnePos(\seq{A189640})(n) = \seq{A026232}(n)$ for all $n \ge 1$, bypassing finite-state transducers altogether (we also note that the index-to-position map $n \mapsto \seq{A026232}(n)$ is not $3$-synchronized). The entire proof is formalized and machine-checked in Lean~4 (\texttt{v4.34.1}, Mathlib \texttt{d13f23b7}) with zero \texttt{sorry}s, allowlist-only axioms, and fresh kernel replay.
\end{abstract}

\section{Introduction, Primary Sources, and Prior-Art Boundary}

\subsection{The Morphism Fixed Point \seq{A189640} and Kimberling's Positional Conjectures}

Consider the $3$-uniform morphism $\alpha : \{0, 1\}^* \to \{0, 1\}^*$ defined on letters by
\begin{equation}
\label{eq:morph}
\alpha(0) = 001, \qquad \alpha(1) = 101.
\end{equation}
Because $\alpha(0)$ begins with $0$ and $\alpha(1)$ begins with $1$, $\alpha$ admits \emph{two} infinite fixed points in $\{0, 1\}^{\N_{\ge 1}}$: one generated from seed $0$ and one from seed $1$, differing solely at index~$1$. In OEIS~\seq{A189640}~\cite{oeis189640} (authored by Clark Kimberling on April~24, 2011; live revision \#45, January~22, 2026; offset~$1$), the seed is fixed unambiguously to $0$ (\texttt{Start: 0}, initial word $W_0 = 0$), yielding the sequence of iterates
\begin{align*}
W_0 &= 0, \\
W_1 &= 001, \\
W_2 &= 001\,001\,101, \\
W_3 &= 001001101\,001001101\,101001101,
\end{align*}
whose limit is the $1$-based binary sequence
\[
\seq{A189640} = (0, 0, 1, 0, 0, 1, 1, 0, 1, 0, 0, 1, 0, 0, 1, 1, 0, 1, 1, 0, 1, 0, 0, 1, 1, 0, 1, \dots).
\]
Its $1$-based $0$-positions begin $1, 2, 4, 5, 8, 10, 11, 13, 14, \dots$, and its $1$-based $1$-positions begin
\[
\OnePos(\seq{A189640}) = (3, 6, 7, 9, 12, 15, 16, 18, 19, 21, 24, 25, 27, 30, 33, \dots).
\]
In the comment block of \seq{A189640} (line~11), Kimberling posed two open positional questions:
\begin{quote}\itshape
Are the positions of $0$ given by \seq{A026138}? Are the positions of $1$ given by \seq{A026232}?
\end{quote}
The first question (designated \texttt{C017}) was proved by F.~M.\ Dekking~\cite{dekking2020} (Theorem~4 and Proposition~6). The second question (designated \texttt{C018}) is the target of the present paper.

\begin{remark}[Conventions and the \seq{A189640} \texttt{\%F} offset slip]
\label{rem:conventions}
All six OEIS sequences in scope (\seq{A189640}, \seq{A026138}, \seq{A026232}, \seq{A026231}, \seq{A026166}, \seq{A026142}) carry OEIS offset~$1$. We maintain $1$-based indexing throughout, so that $\seq{A189640}(1) = 0$, $\seq{A189640}(2) = 0$, $\seq{A189640}(3) = 1$, and $\OnePos(\seq{A189640})(1) = 3$. We note that line~49 of \seq{A189640}'s \texttt{\%F} field prints ``$a(3k-2)=a(k),\; a(3k-1)=0,\; a(3k)=1$ for $k \ge 1$, $a(0)=0$'', which contains a minor base-index slip ($a(0)=0$ in place of $a(1)=0$, leaving $k=1$ tautological as $a(1)=a(1)$). To avoid any circularity, our definition and Lean~4 formalization construct $\seq{A189640}$ directly from the morphism iterates $W_m = \alpha^m(0)$ and derive the $3$-way recurrence unconditionally from~\eqref{eq:morph}.
\end{remark}

\subsection{Left-Right Filling Permutations \seq{A026142}, \seq{A026166} and Derived Sequences \seq{A026231}, \seq{A026232}}

Given a step-offset function $h : \N \to \N$, Kimberling's \emph{left-right filling algorithm} constructs an array indexed by $\N_{\ge 1}$ as follows. Initialize step~$1$ by placing value~$1$ at index~$1$: $a(1) = 1$. Then, successively for each integer $n = 2, 3, 4, \dots$, set
\begin{equation}
\label{eq:lr-rule}
L_h(n) := n - h(n), \qquad R_h(n) := n + h(n),
\end{equation}
and place $n$ into cell $L_h(n)$ if that cell is currently empty, and into cell $R_h(n)$ otherwise.

The two permutations governing \seq{A026231} and \seq{A026232} arise from the two adjacent half-integer step offsets:
\begin{itemize}
\item $\Pi_{66} := \seq{A026166}$~\cite{oeis026166} (live revision \#35, July~3, 2020): $h_{66}(n) = \lfloor (n - 1)/2 \rfloor$, with initial condition $\Pi_{66}(1) = 1$. (Although the \texttt{\%N} title sentence of \seq{A026166} begins ``For $n \ge 2$, \dots'', the \texttt{\%S} data, Mathematica/Python programs, jOEIS implementation, and Dekking~\cite[§1]{dekking2020} all initialize $\Pi_{66}(1) = 1$; indeed $L_{66}(n) = \lfloor (n+2)/2 \rfloor \ge 2$ for every $n \ge 2$, so cell~$1$ is never reached at $n \ge 2$.)
\item $\Pi_{42} := \seq{A026142}$~\cite{oeis026142} (live revision \#16, January~14, 2022): $h_{42}(n) = \lfloor (n + 1)/2 \rfloor$, with initial condition $\Pi_{42}(1) = 1$. On January~7, 2022, Peter Munn renamed \seq{A026142} to the ``lexicographically earliest permutation of the positive integers such that for $n \ge 2$, $|a(n) - n| = \lfloor (a(n)+1)/2 \rfloor$'', moving Kimberling's left-right filling specification to \texttt{\%C} (\emph{Old name}). As verified directly from the residue-class description in Section~\ref{sec:residues} (and Munn's comment on \seq{A026142}), $\Pi_{42}$ achieves the pointwise minimum unused candidate at every index $j \ge 1$, so both definitions determine the exact same sequence.
\end{itemize}
From $\Pi_{66}$ and $\Pi_{42}$, Kimberling defines~\cite{oeis026231,oeis026232}:
\begin{align}
\seq{A026231}(n) &:= \text{the $n$-th positive integer $K \ge 1$ such that } \Pi_{66}(K) = \Pi_{42}(K) + 1, \label{eq:def231} \\
\seq{A026232}(n) &:= \frac{\seq{A026231}(n) + 1}{3} \qquad (n \ge 1). \label{eq:def232}
\end{align}
Table~\ref{tab:small-k} records $\Pi_{66}(k)$, $\Pi_{42}(k)$, and their differences across $k = 1, \dots, 30$. The first four coincidence indices $\Pi_{66}(k) = \Pi_{42}(k) + 1$ are $k \in \{8, 17, 20, 26\}$, all congruent to $2 \pmod{3}$, giving $(k+1)/3 \in \{3, 6, 7, 9\} = \OnePos(\seq{A189640})(1..4)$.

\begin{table}[ht]
\centering
\caption{Values of $\Pi_{66}(k) = \seq{A026166}(k)$, $\Pi_{42}(k) = \seq{A026142}(k)$, and $\Pi_{66}(k) - \Pi_{42}(k)$ for $k = 1, \dots, 30$. Bold rows mark the $+1$ coincidences \seq{A026231} and their images $(k+1)/3 = \seq{A026232}$.}
\label{tab:small-k}
\smallskip
{\small
\begin{tabular}{@{}rrrrcc|rrrrcc@{}}
\toprule
$k$ & $\Pi_{66}(k)$ & $\Pi_{42}(k)$ & Diff & \seq{A026231} & \seq{A026232} & $k$ & $\Pi_{66}(k)$ & $\Pi_{42}(k)$ & Diff & \seq{A026231} & \seq{A026232} \\
\midrule
1  & 1  & 1  & $0$   & --- & --- & 16 & 11 & 32 & $-21$ & --- & --- \\
2  & 2  & 4  & $-2$  & --- & --- & \textbf{17} & \textbf{12} & \textbf{11} & $\mathbf{+1}$ & \textbf{17} & \textbf{6} \\
3  & 4  & 2  & $+2$  & --- & --- & 18 & 34 & 36 & $-2$  & --- & --- \\
4  & 3  & 8  & $-5$  & --- & --- & 19 & 13 & 38 & $-25$ & --- & --- \\
5  & 8  & 3  & $+5$  & --- & --- & \textbf{20} & \textbf{14} & \textbf{13} & $\mathbf{+1}$ & \textbf{20} & \textbf{7} \\
6  & 10 & 12 & $-2$  & --- & --- & 21 & 40 & 42 & $-2$  & --- & --- \\
7  & 5  & 14 & $-9$  & --- & --- & 22 & 15 & 44 & $-29$ & --- & --- \\
\textbf{8} & \textbf{6} & \textbf{5} & $\mathbf{+1}$ & \textbf{8} & \textbf{3} & 23 & 44 & 15 & $+29$ & --- & --- \\
9  & 16 & 6  & $+10$ & --- & --- & 24 & 46 & 16 & $+30$ & --- & --- \\
10 & 7  & 20 & $-13$ & --- & --- & 25 & 17 & 50 & $-33$ & --- & --- \\
11 & 20 & 7  & $+13$ & --- & --- & \textbf{26} & \textbf{18} & \textbf{17} & $\mathbf{+1}$ & \textbf{26} & \textbf{9} \\
12 & 22 & 24 & $-2$  & --- & --- & 27 & 52 & 18 & $+34$ & --- & --- \\
13 & 9  & 26 & $-17$ & --- & --- & 28 & 19 & 56 & $-37$ & --- & --- \\
14 & 26 & 9  & $+17$ & --- & --- & 29 & 56 & 19 & $+37$ & --- & --- \\
15 & 28 & 10 & $+18$ & --- & --- & 30 & 58 & 60 & $-2$  & --- & --- \\
\bottomrule
\end{tabular}
}
\end{table}

\subsection{Prior-Art Boundary}

The closest published work is F.~M.\ Dekking's paper \emph{Permutations of $\N$ generated by left-right filling algorithms}, posted as \texttt{arXiv:2001.08915v1} (January~24, 2020) and published in the \emph{Journal of Integer Sequences}, Vol.~23 (2020), Article~20.4.7~\cite{dekking2020}. First-hand inspection of the published \LaTeX{} source (\texttt{dekk5.tex}) confirms the exact boundary of prior art:
\begin{enumerate}
\item Dekking proved the companion $0$-positions conjecture \texttt{C017} ($\ZeroPos(\seq{A189640}) = \seq{A026138}$, Theorem~4 and Proposition~6) and analyzed \seq{A026136}, \seq{A026142} (Section~2.4, Theorems~9--10), \seq{A026172}, and \seq{A026177}.
\item In Section~3 (\emph{Conclusion}, line~570 of \texttt{dekk5.tex}), Dekking explicitly left \seq{A026166} open:
\begin{quote}\itshape
This has lead to a compact description of sequences \seq{A026136}, \seq{A026142}, \seq{A026172}, and \seq{A026177} in the OEIS. We leave the permutation given in \seq{A026166} to the interested reader.
\end{quote}
Neither \seq{A026231} nor \seq{A026232} is analyzed in the body of~\cite{dekking2020} (they appear only in the journal's footer list of related OEIS entries), and \seq{A189640} does not appear anywhere in~\cite{dekking2020}.
\item Dekking's October~26, 2019 comment on \seq{A189640} proves Mathar's shift identity $\seq{A189640}(n + 1) = \seq{A080846}(n)$ for $n \ge 1$~\cite{oeis080846}, which does not resolve the positional identity with \seq{A026232}.
\item Live re-verification on November~24, 2026 across live web searches, live OEIS records (\seq{A189640} \#45, \seq{A026232} \#6, \seq{A026231} \#14, \seq{A026166} \#35, \seq{A026142} \#16), the arXiv API, and OpenAlex confirms that no published proof or preprint resolution of \texttt{C018} exists. (Note: \seq{A189640} carries a batch-added citation to Table~5.1, p.~27 of Mazen Khodier's 2026 University of Waterloo Master's thesis, shared across Kimberling's binary $3$-uniform morphisms of \seq{A189628} for subword-complexity benchmarks; UWSpace is gated by a proof-of-work challenge that was not accessed or bypassed, and none of \seq{A026232}, \seq{A026231}, \seq{A026166}, or \seq{A026142} cites Khodier, while \seq{A189640} retains \texttt{C018} as an open question.)
\end{enumerate}

\subsection{Statement of Main Results}

\begin{theorem}[Complete residue-class characterization of \seq{A026166}]
\label{thm:a026166}
The left-right filling sequence $\Pi_{66} = \seq{A026166}$ has no cell overwrites, satisfies $\Pi_{66}(1) = 1$, $\Pi_{66}(2) = 2$, and for every integer $r \ge 1$:
\begin{align}
\Pi_{66}(3r) &= 6r - 2, \label{eq:66-0} \\
\Pi_{66}(3r + 1) &= 2r + 1, \label{eq:66-1} \\
\Pi_{66}(3r + 2) &= \begin{cases}
6r + 2 & \text{if } \Pi_{66}(r + 2) = 2r + 2, \\
2r + 2 & \text{if } \Pi_{66}(r + 2) \ne 2r + 2.
\end{cases} \label{eq:66-2}
\end{align}
Moreover, for every $t \in \N$, $\Pi_{66}(t + 2) \ne 2t + 2$ if and only if $\seq{A189640}(t + 1) = 1$.
\end{theorem}

\begin{theorem}[Exact division-by-$3$ integrality and positional identity \texttt{C018}]
\label{thm:main}
For every integer $n \ge 1$:
\begin{enumerate}
\item[\textup{(a)}] $\seq{A026231}(n) = 3\,\OnePos(\seq{A189640})(n) - 1$;
\item[\textup{(b)}] $\seq{A026231}(n) \equiv 2 \pmod{3}$, so that $\frac{1}{3}(\seq{A026231}(n) + 1) = \lfloor (\seq{A026231}(n) + 1)/3 \rfloor$ is an exact integer quotient with zero remainder;
\item[\textup{(c)}] $\OnePos(\seq{A189640})(n) = \seq{A026232}(n)$.
\end{enumerate}
\end{theorem}

\section{Morphism Fixed Point and Ternary Recurrence for \seq{A189640}}
\label{sec:morph}

Define $W_0 := [0] \in \{0, 1\}^*$ and $W_{m+1} := \alpha(W_m)$ for $m \in \N$, where $\alpha(w)$ applies~\eqref{eq:morph} letter-by-letter (\texttt{List.flatMap}). Because $|\alpha(b)| = 3$ for both $b \in \{0, 1\}$, induction on $|w|$ gives $|\alpha(w)| = 3|w|$, and induction on $m$ gives $|W_m| = 3^m$.

\begin{lemma}[Prefix stability and block indexing]
\label{lem:morph-block}
\begin{enumerate}
\item[\textup{(a)}] For every word $w \in \{0, 1\}^*$ and every $0$-based index $j < |w|$,
\[
\alpha(w)_{3j} = w_j, \qquad \alpha(w)_{3j+1} = 0, \qquad \alpha(w)_{3j+2} = 1.
\]
\item[\textup{(b)}] For all $m_1 \le m_2$, $W_{m_1}$ is a prefix of $W_{m_2}$; consequently, if $j < 3^{m_1}$, then $(W_{m_1})_j = (W_{m_2})_j$.
\end{enumerate}
\end{lemma}

\begin{proof}
Part~(a) is immediate by induction on $w$: the empty case is vacuous, and for $w = b \cdot u$, $\alpha(b \cdot u) = \alpha(b) \cdot \alpha(u)$ begins with the length-$3$ block $(b, 0, 1)$ for both $b \in \{0, 1\}$, handling $j = 0$, while $j = k + 1$ shifts all three indices by~$3$ into $\alpha(u)$. For part~(b), $\alpha$ preserves the prefix relation ($u \sqsubseteq v \implies \alpha(u) \sqsubseteq \alpha(v)$) since $\alpha(u \cdot r) = \alpha(u) \cdot \alpha(r)$. Because $W_0 = [0] \sqsubseteq [0, 0, 1] = W_1$, induction on $m$ gives $W_m \sqsubseteq W_{m+1}$, and transitivity yields $W_{m_1} \sqsubseteq W_{m_2}$ for all $m_1 \le m_2$.
\end{proof}

Since $n < 3^n$ for all $n \in \N$ (by induction: $0 < 1$ and $n + 1 < 3^n + 1 \le 3^{n+1}$), every $0$-based index $j \ge 0$ satisfies $j < j + 1 < 3^{j+1} = |W_{j+1}|$. Thus the $1$-based sequence
\[
a(n) := (W_n)_{n-1} \qquad (n \ge 1)
\]
is well-defined, and by Lemma~\ref{lem:morph-block}(b) satisfies $a(j + 1) = (W_m)_j$ for every $m$ with $j < 3^m$.

\begin{proposition}[Ternary recurrence and infinitude of $1$-positions of \seq{A189640}]
\label{prop:a189640-rec}
The $1$-based sequence $a = \seq{A189640}$ satisfies $a(1) = 0$ and, for every $j \in \N$:
\begin{equation}
\label{eq:a189640-rec}
a(3j + 1) = a(j + 1), \qquad a(3j + 2) = 0, \qquad a(3j + 3) = 1.
\end{equation}
In particular, the set $\{k \ge 1 : a(k) = 1\}$ is infinite.
\end{proposition}

\begin{proof}
At $n = 1$, $a(1) = (W_1)_0 = 0$. Fix $j \in \N$ and put $m = j + 1$. Then $j < 3^m = |W_m|$, whence $3j < 3j + 1 < 3j + 2 < 3^{m+1}$. By Lemma~\ref{lem:morph-block}(b) and~(a) applied to $W_{m+1} = \alpha(W_m)$,
\begin{align*}
a(3j + 1) &= (W_{m+1})_{3j} = (W_m)_j = a(j + 1), \\
a(3j + 2) &= (W_{m+1})_{3j+1} = 0, \\
a(3j + 3) &= (W_{m+1})_{3j+2} = 1.
\end{align*}
Since $3n + 3 > n$ and $a(3n + 3) = 1$ for every $n \in \N$, $\{k \ge 1 : a(k) = 1\}$ is unbounded and therefore infinite.
\end{proof}

\section{Uniform No-Overwrite Persistence for Left-Right Filling}
\label{sec:lr}

For any step-offset function $h : \N \to \N$, we formalize the partial array state $S_h : \N \to \N \to \mathrm{Option}\,\N$ recursively by $S_h(0, k) = \bot$, $S_h(1, k) = \mathrm{some}\,1$ if $k = 1$ else $\bot$, and for $n \ge 2$,
\[
S_h(n, k) := \begin{cases}
\mathrm{some}\,n & \text{if } k = T_h(n), \\
S_h(n - 1, k) & \text{otherwise},
\end{cases}
\qquad
T_h(n) := \begin{cases}
L_h(n) & \text{if } S_h(n - 1, L_h(n)) = \bot, \\
R_h(n) & \text{otherwise},
\end{cases}
\]
where $L_h(n) = n - h(n)$ and $R_h(n) = n + h(n)$. The total sequence $\mathrm{lrFill}_h(k)$ reads the value stored at cell~$k$ at the least step $m = \inf\{m' \in \N : S_h(m', k) \ne \bot\}$.

\begin{definition}[Regular left-right offset]
\label{def:good-lr}
A step-offset function $h : \N \to \N$ is \emph{regular} (\texttt{GoodLR h}) if:
\begin{enumerate}
\item[\textup{(i)}] $1 < R_h(n)$ for all $n \ge 2$;
\item[\textup{(ii)}] $L_h(v) < R_h(n)$ for all $2 \le v < n$;
\item[\textup{(iii)}] $R_h(v) < R_h(n)$ for all $2 \le v < n$ (strict monotonicity of $R_h$ on $\N_{\ge 2}$).
\end{enumerate}
\end{definition}

Note that both $h_{42}(n) = \lfloor (n+1)/2 \rfloor$ and $h_{66}(n) = \lfloor (n-1)/2 \rfloor$ are regular: for $n \ge 2$, $R_{42}(n) \ge 3 > 1$ and $R_{66}(n) \ge 2 > 1$; and for $2 \le v < n$ (forcing $n \ge 3$),
\begin{align*}
L_{42}(v) &= \lfloor v/2 \rfloor < n < R_{42}(n), \\
R_{42}(v) &= v + \lfloor (v+1)/2 \rfloor < n + \lfloor (n+1)/2 \rfloor = R_{42}(n), \\
L_{66}(v) &= \lfloor (v+2)/2 \rfloor \le \lfloor (n+1)/2 \rfloor < n + \lfloor (n-1)/2 \rfloor = R_{66}(n), \\
R_{66}(v) &= v + \lfloor (v-1)/2 \rfloor < n + \lfloor (n-1)/2 \rfloor = R_{66}(n).
\end{align*}

\begin{lemma}[No-overwrite persistence and exact extraction]
\label{lem:lr-persist}
Let $h : \N \to \N$ be any step-offset function.
\begin{enumerate}
\item[\textup{(a)}] \textup{(Support bound.)} Whenever $S_h(m, k) = \mathrm{some}\,v$, either $(k = 1 \wedge v = 1 \wedge m \ge 1)$ or $(2 \le v \le m \wedge k \in \{L_h(v), R_h(v)\})$. In particular, $v \le m$.
\item[\textup{(b)}] \textup{(Step-range emptiness preservation.)} If $1 \le m_1 \le m_2$, $S_h(m_1, k) = \bot$, and $k \notin \{L_h(v), R_h(v)\}$ for every $v$ with $m_1 < v \le m_2$, then $S_h(m_2, k) = \bot$.
\end{enumerate}
Assume further that $h$ is regular (Definition~\ref{def:good-lr}). Then:
\begin{enumerate}
\item[\textup{(c)}] \textup{(Freshness of targets.)} For every $n \ge 2$, $S_h(n - 1, T_h(n)) = \bot$.
\item[\textup{(d)}] \textup{(Persistence and uniqueness.)} If $m_1 \le m_2$ and $S_h(m_1, k) = \mathrm{some}\,v$, then $S_h(m_2, k) = \mathrm{some}\,v$. Consequently, if $S_h(m, k) = \mathrm{some}\,v$ for any step~$m$, then $\mathrm{lrFill}_h(k) = v$.
\item[\textup{(e)}] \textup{(Guaranteed occupation of $L_h(n)$.)} For every $n \ge 2$, $S_h(n, L_h(n)) \ne \bot$.
\end{enumerate}
\end{lemma}

\begin{proof}
Part~(a) follows by straightforward induction on $m$, since step~$1$ writes only $(1 \mapsto 1)$ and step~$n \ge 2$ writes only $(T_h(n) \mapsto n)$ with $T_h(n) \in \{L_h(n), R_h(n)\}$. Part~(b) follows by induction on $m_2 \ge m_1$, since at each step $v \in (m_1, m_2]$ we have $v \ge 2$ and $T_h(v) \in \{L_h(v), R_h(v)\}$, so $k \ne T_h(v)$ leaves cell~$k$ untouched.

Now suppose $h$ is regular. For~(c), if $S_h(n-1, L_h(n)) = \bot$ then $T_h(n) = L_h(n)$ and the claim holds trivially. Otherwise $T_h(n) = R_h(n)$; if $S_h(n-1, R_h(n)) = \mathrm{some}\,v$, then part~(a) gives either $R_h(n) = 1$ (contradicting~(i)) or $2 \le v \le n - 1$ with $R_h(n) \in \{L_h(v), R_h(v)\}$ (contradicting~(ii) or~(iii)). Thus $S_h(n-1, T_h(n)) = \bot$.

For~(d), induction on $m_2 \ge m_1$ shows persistence: at step $m_2 = m + 2$, if $k = T_h(m+2)$ then $S_h(m+1, k) = \bot$ by~(c), contradicting the inductive hypothesis $S_h(m+1, k) = \mathrm{some}\,v$; hence $k \ne T_h(m+2)$ and $S_h(m+2, k) = \mathrm{some}\,v$. Any two steps $m_1, m_2$ with $S_h(m_i, k) = \mathrm{some}\,v_i$ therefore agree ($v_1 = v_2$) by lifting the smaller step to $\max(m_1, m_2)$. In particular, if $S_h(m, k) = \mathrm{some}\,v$, then $\{m' : S_h(m', k) \ne \bot\}$ is nonempty, so its minimum $m_0$ satisfies $S_h(m_0, k) = \mathrm{some}\,w$ for some $w$, forcing $w = v$ and hence $\mathrm{lrFill}_h(k) = v$.

For~(e), if $S_h(n-1, L_h(n)) = \bot$, step~$n$ writes $\mathrm{some}\,n$ into $L_h(n)$; otherwise $S_h(n-1, L_h(n)) = \mathrm{some}\,v$ persists to step~$n$ by~(d).
\end{proof}

\section{Residue Classes of \seq{A026142} and Complete Characterization of \seq{A026166}}
\label{sec:residues}

\subsection{Residue Classes of \seq{A026142}}

For $h_{42}(n) = \lfloor (n+1)/2 \rfloor$, evaluating the left-right endpoints on parity classes ($j \ge 1$) yields
\begin{equation}
\label{eq:LR42}
L_{42}(2j) = j,\quad R_{42}(2j) = 3j, \qquad
L_{42}(2j + 1) = j,\quad R_{42}(2j + 1) = 3j + 2.
\end{equation}
Direct execution of steps $1, \dots, 8$ together with Lemma~\ref{lem:lr-persist}(d) gives the small values $\Pi_{42}(1) = 1$, $\Pi_{42}(2) = 4$, $\Pi_{42}(3) = 2$, and $\Pi_{42}(4) = 8$.

\begin{proposition}[Residue-class formulas for \seq{A026142}]
\label{prop:a026142}
For every integer $r \ge 1$:
\begin{enumerate}
\item[\textup{(a)}] $\Pi_{42}(3r + 2) = 2r + 1$;
\item[\textup{(b)}] $\Pi_{42}(3r + 1) = 6r + 2$;
\item[\textup{(c)}] $\Pi_{42}(3r) \in \{2r, 6r\}$ (in particular, $\Pi_{42}(3r)$ is always even).
\end{enumerate}
\end{proposition}

\begin{proof}
(a) At step $2r \ge 2$, $L_{42}(2r) = r$, so cell~$r$ is occupied after step~$2r$ by Lemma~\ref{lem:lr-persist}(e). Therefore step $2r + 1$ (which has $L_{42}(2r + 1) = r$ and $R_{42}(2r + 1) = 3r + 2$) takes the $R$-branch and writes $S_{h_{42}}(2r + 1, 3r + 2) = \mathrm{some}\,(2r + 1)$, whence $\Pi_{42}(3r + 2) = 2r + 1$ by Lemma~\ref{lem:lr-persist}(d).

(b) By~\eqref{eq:LR42}, $R_{42}(v) \in \{0, 2\} \pmod{3}$ for every $v \ge 2$, so $R_{42}(v) \ne 3r + 1$. Moreover, for all $2 \le v \le 6r + 1$, $L_{42}(v) = \lfloor v/2 \rfloor \le 3r < 3r + 1$, and $3r + 1 \ne 1$. By Lemma~\ref{lem:lr-persist}(a), $S_{h_{42}}(6r + 1, 3r + 1) = \bot$. Since $L_{42}(6r + 2) = 3r + 1$, step $6r + 2$ takes the $L$-branch and sets $\Pi_{42}(3r + 1) = 6r + 2$.

(c) Inspect $S_{h_{42}}(6r - 1, 3r)$. If it equals $\mathrm{some}\,v$, then Lemma~\ref{lem:lr-persist}(a) gives $2 \le v \le 6r - 1$ (since $3r \ne 1$) with $3r \in \{L_{42}(v), R_{42}(v)\}$. Because $L_{42}(v) = \lfloor v/2 \rfloor \le 3r - 1 < 3r$, we must have $R_{42}(v) = 3r$, which by~\eqref{eq:LR42} forces $v = 2r$, giving $\Pi_{42}(3r) = 2r$. Otherwise $S_{h_{42}}(6r - 1, 3r) = \bot$; since $L_{42}(6r) = 3r$, step~$6r$ fills cell~$3r$ with $6r$, giving $\Pi_{42}(3r) = 6r$.
\end{proof}

\begin{remark}
In Dekking~\cite[Theorem~9]{dekking2020}, parts~(a) and~(b) are stated for $k \in \N_{\ge 1}$ (note that at $r = 0$, $\Pi_{42}(1) = 1 \ne 2$ and $\Pi_{42}(2) = 4 \ne 1$). Notice also that for \texttt{C018} we do not even need the deeper recursive split of $\Pi_{42}(3r)$ between $2r$ and $6r$—the parity consequence $\Pi_{42}(3r) \in \{2r, 6r\}$ suffices completely.
\end{remark}

\subsection{Complete Characterization of \seq{A026166} (Proof of Theorem~\ref{thm:a026166})}

For $h_{66}(n) = \lfloor (n-1)/2 \rfloor$, the left-right endpoints on parity classes ($j \ge 1$) are
\begin{equation}
\label{eq:LR66}
L_{66}(2j) = j + 1,\quad R_{66}(2j) = 3j - 1, \qquad
L_{66}(2j + 1) = j + 1,\quad R_{66}(2j + 1) = 3j + 1.
\end{equation}
Direct execution of steps $1, \dots, 4$ with Lemma~\ref{lem:lr-persist}(d) gives $\Pi_{66}(1) = 1$, $\Pi_{66}(2) = 2$, $\Pi_{66}(3) = 4$, and $\Pi_{66}(4) = 3$.

\begin{proof}[Proof of~\eqref{eq:66-0}--\eqref{eq:66-2} in Theorem~\ref{thm:a026166}]
Fix $r \ge 1$.

\emph{Case $3r$ (Equation~\eqref{eq:66-0}):} By~\eqref{eq:LR66}, $R_{66}(v) \in \{1, 2\} \pmod{3}$ for every $v \ge 2$, so $R_{66}(v) \ne 3r$. For all $2 \le v \le 6r - 3$, $L_{66}(v) = \lfloor (v + 2)/2 \rfloor \le 3r - 1 < 3r$, and $3r \ne 1$. By Lemma~\ref{lem:lr-persist}(a), $S_{h_{66}}(6r - 3, 3r) = \bot$. Since $L_{66}(6r - 2) = 3r$ (with $6r - 2 \ge 4$), step $6r - 2$ writes $\Pi_{66}(3r) = 6r - 2$.

\emph{Case $3r + 1$ (Equation~\eqref{eq:66-1}):} At step $2r \ge 2$, $L_{66}(2r) = r + 1$, so cell $r + 1$ is occupied after step~$2r$ by Lemma~\ref{lem:lr-persist}(e). Hence step $2r + 1$ (with $L_{66}(2r+1) = r + 1$ and $R_{66}(2r+1) = 3r + 1$) takes the $R$-branch and writes $\Pi_{66}(3r + 1) = 2r + 1$.

\emph{Case $3r + 2$ (Equation~\eqref{eq:66-2}, the key branching lemma):} Consider step $2r + 2$ ($j = r + 1 \ge 2$), which has $L_{66}(2r + 2) = r + 2$ and $R_{66}(2r + 2) = 3r + 2$. Split on the state of cell $r + 2$ immediately before step $2r + 2$:
\begin{itemize}
\item \emph{Subcase 1: $S_{h_{66}}(2r + 1, r + 2) = \bot$.} Then step $2r + 2$ takes the $L$-branch, writing $S_{h_{66}}(2r + 2, r + 2) = \mathrm{some}\,(2r + 2)$, so $\Pi_{66}(r + 2) = 2r + 2$ by Lemma~\ref{lem:lr-persist}(d). Moreover, since $r \ge 1$, the target $r + 2$ is distinct from $3r + 2$. Before step $2r + 2$, every $2 \le v \le 2r + 1$ satisfies $L_{66}(v) \le r + 1 < 3r + 2$ and $R_{66}(v) \le 3r + 1 < 3r + 2$, so Lemma~\ref{lem:lr-persist}(a) gives $S_{h_{66}}(2r + 1, 3r + 2) = \bot$; since step $2r + 2$ writes to $r + 2 \ne 3r + 2$, we also have $S_{h_{66}}(2r + 2, 3r + 2) = \bot$.

Now consider every intermediate step $v$ with $2r + 2 < v \le 6r + 1$:
\[
L_{66}(v) = \lfloor (v + 2)/2 \rfloor \le 3r + 1 < 3r + 2,
\]
and by strict monotonicity of $R_{66}$,
\[
R_{66}(v) \ge R_{66}(2r + 3) = 3r + 4 > 3r + 2.
\]
Hence cell $3r + 2$ is never targeted during steps $2r + 3, \dots, 6r + 1$, so Lemma~\ref{lem:lr-persist}(b) yields $S_{h_{66}}(6r + 1, 3r + 2) = \bot$. Finally, $L_{66}(6r + 2) = 3r + 2$, so step $6r + 2$ fills cell $3r + 2$ via the $L$-branch, giving $\Pi_{66}(3r + 2) = 6r + 2$.

\item \emph{Subcase 2: $S_{h_{66}}(2r + 1, r + 2) = \mathrm{some}\,v$.} By Lemma~\ref{lem:lr-persist}(a) and~(d), $\Pi_{66}(r + 2) = v \le 2r + 1$, so $\Pi_{66}(r + 2) \ne 2r + 2$. Because cell $r + 2$ is already occupied before step $2r + 2$, step $2r + 2$ takes the $R$-branch and writes $S_{h_{66}}(2r + 2, 3r + 2) = \mathrm{some}\,(2r + 2)$, giving $\Pi_{66}(3r + 2) = 2r + 2$.
\end{itemize}
Combining both subcases establishes~\eqref{eq:66-2}.
\end{proof}

\section{Recurrence Coincidence, Set Equivalence, and Lifting to Sequence Equality}
\label{sec:main}

\subsection{Coincidence of the Two Ternary Recurrences}

We now complete Theorem~\ref{thm:a026166} by identifying the branching condition of~\eqref{eq:66-2} directly with $\seq{A189640}$.

\begin{proposition}[Ternary recurrence match]
\label{prop:rec-match}
For every $t \in \N$,
\[
\Pi_{66}(t + 2) \ne 2t + 2 \iff \seq{A189640}(t + 1) = 1.
\]
\end{proposition}

\begin{proof}
We proceed by strong induction on $t \in \N$.
\begin{itemize}
\item For $t = 0$: $\Pi_{66}(2) = 2 = 2(0) + 2$ and $\seq{A189640}(1) = 0 \ne 1$, so both sides are false.
\item For $t \ge 1$, write $t = 3u + c$ with $u = \lfloor t/3 \rfloor$ and $c \in \{0, 1, 2\}$:
\begin{itemize}
\item If $c = 0$, then $t = 3u$ with $1 \le u < t$. By~\eqref{eq:66-2}, since $2u + 2 < 6u + 2$,
\[
\Pi_{66}(3u + 2) \ne 6u + 2 \iff \Pi_{66}(u + 2) \ne 2u + 2.
\]
By the inductive hypothesis at $u < t$ and Proposition~\ref{prop:a189640-rec},
\[
\Pi_{66}(u + 2) \ne 2u + 2 \iff \seq{A189640}(u + 1) = 1 \iff \seq{A189640}(3u + 1) = 1.
\]
\item If $c = 1$, then $t + 2 = 3(u + 1)$. Equation~\eqref{eq:66-0} gives $\Pi_{66}(3(u+1)) = 6(u+1) - 2 = 6u + 4 = 2(3u + 1) + 2$, while Proposition~\ref{prop:a189640-rec} gives $\seq{A189640}(3u + 2) = 0 \ne 1$; both sides are false.
\item If $c = 2$, then $t + 2 = 3(u + 1) + 1$. Equation~\eqref{eq:66-1} gives $\Pi_{66}(3(u+1) + 1) = 2(u+1) + 1 = 2u + 3 \ne 6u + 6 = 2(3u + 2) + 2$, while Proposition~\ref{prop:a189640-rec} gives $\seq{A189640}(3u + 3) = 1$; both sides are true. \qedhere
\end{itemize}
\end{itemize}
\end{proof}

\subsection{Set-Level Characterization of \seq{A026231}}

\begin{proposition}[Set equivalence]
\label{prop:set-equiv}
For every $K \in \N$,
\[
\bigl(1 \le K \;\wedge\; \Pi_{66}(K) = \Pi_{42}(K) + 1\bigr) \iff \exists\, m \in \N,\; \bigl(1 \le m \;\wedge\; \seq{A189640}(m) = 1\bigr) \;\wedge\; K = 3m - 1.
\]
\end{proposition}

\begin{proof}
$(\Rightarrow)$ For $K \in \{1, 2, 3, 4\}$, Table~\ref{tab:small-k} gives $\Pi_{66}(K) - \Pi_{42}(K) \in \{0, -2, +2, -5\}$, none equal to $+1$. For $K \ge 5$, write $K = 3r + c$ with $r = \lfloor K/3 \rfloor \ge 1$ and $c \in \{0, 1, 2\}$:
\begin{itemize}
\item If $c = 0$, then $r \ge 2$, $\Pi_{66}(3r) = 6r - 2$ is even by~\eqref{eq:66-0}, whereas $\Pi_{42}(3r) + 1 \in \{2r + 1, 6r + 1\}$ is odd by Proposition~\ref{prop:a026142}(c)—impossible.
\item If $c = 1$, then $r \ge 2$, and~\eqref{eq:66-1} versus Proposition~\ref{prop:a026142}(b) gives $\Pi_{66}(3r + 1) = 2r + 1 < 6r + 3 = \Pi_{42}(3r + 1) + 1$—impossible.
\item If $c = 2$, then $r \ge 1$, and Proposition~\ref{prop:a026142}(a) gives $\Pi_{42}(3r + 2) + 1 = 2r + 2$. Since $6r + 2 > 2r + 2$,~\eqref{eq:66-2} forces $\Pi_{66}(r + 2) \ne 2r + 2$. By Proposition~\ref{prop:rec-match}, $\seq{A189640}(r + 1) = 1$. Setting $m := r + 1 \ge 2$ satisfies $1 \le m$, $\seq{A189640}(m) = 1$, and $K = 3r + 2 = 3m - 1$.
\end{itemize}

$(\Leftarrow)$ Given $m \ge 1$ with $\seq{A189640}(m) = 1$ and $K = 3m - 1$, we have $m \ge 2$ because $\seq{A189640}(1) = 0$. Write $m = r + 1$ ($r \ge 1$), so $K = 3r + 2 \ge 5 \ge 1$. Proposition~\ref{prop:rec-match} gives $\Pi_{66}(r + 2) \ne 2r + 2$, hence~\eqref{eq:66-2} gives $\Pi_{66}(3r + 2) = 2r + 2 = (2r + 1) + 1 = \Pi_{42}(3r + 2) + 1$ by Proposition~\ref{prop:a026142}(a).
\end{proof}

\subsection{Order-Preserving Enumeration Lift (Proof of Theorem~\ref{thm:main})}

\begin{lemma}[Uniqueness of strictly monotone enumerations]
\label{lem:strict-mono-eq}
If $f, g : \N \to \N$ are strictly increasing functions with the same image $\{f(i) : i \in \N\} = \{g(i) : i \in \N\}$, then $f(i) = g(i)$ for all $i \in \N$.
\end{lemma}

\begin{proof}
By strong induction on $i \in \N$: assuming $f(j) = g(j)$ for all $j < i$, pick $a \in \N$ with $f(a) = g(i)$. If $a < i$, then $g(i) = f(a) = g(a) < g(i)$, impossible; hence $i \le a$ and $f(i) \le f(a) = g(i)$. Symmetrically $g(i) \le f(i)$, so $f(i) = g(i)$.
\end{proof}

\begin{proof}[Proof of Theorem~\ref{thm:main}]
Let $P(m) := (1 \le m \wedge \seq{A189640}(m) = 1)$ and $Q(K) := (1 \le K \wedge \Pi_{66}(K) = \Pi_{42}(K) + 1)$. By Proposition~\ref{prop:a189640-rec}, $\{m : P(m)\}$ is infinite, so its $0$-based increasing enumeration $i \mapsto \mathrm{nth}_P(i)$ is strictly increasing with image $\{m : P(m)\} \subseteq \N_{\ge 1}$. By Proposition~\ref{prop:set-equiv}, $Q(K) \iff \exists\, m,\, P(m) \wedge K = 3m - 1$, so $\{K : Q(K)\}$ is also infinite and its increasing enumeration $f(i) := \mathrm{nth}_Q(i)$ is strictly increasing with image $\{K : Q(K)\}$.

Define $g(i) := 3\,\mathrm{nth}_P(i) - 1$. Since $1 \le \mathrm{nth}_P(i) < \mathrm{nth}_P(j)$ implies $3\,\mathrm{nth}_P(i) - 1 < 3\,\mathrm{nth}_P(j) - 1$, $g$ is strictly increasing, and by Proposition~\ref{prop:set-equiv} the image of $g$ is exactly $\{K : Q(K)\}$. Lemma~\ref{lem:strict-mono-eq} therefore yields $\mathrm{nth}_Q(i) = 3\,\mathrm{nth}_P(i) - 1$ for every $i \in \N$. Setting $i = n - 1$ for $n \ge 1$ gives part~(a):
\[
\seq{A026231}(n) = 3\,\OnePos(\seq{A189640})(n) - 1.
\]
Since $\OnePos(\seq{A189640})(n) \ge 1$, writing $M := \OnePos(\seq{A189640})(n) \ge 1$ yields $(3M - 1) \bmod 3 = 2$ (part~(b)) and
\[
\seq{A026232}(n) = \frac{(3M - 1) + 1}{3} = \frac{3M}{3} = M = \OnePos(\seq{A189640})(n),
\]
proving part~(c).
\end{proof}

\subsection{Remark on Automata, Non-Synchronizability, and Avoidance of Walnut}

\begin{remark}[3-automatic characteristic sets vs.\ non-synchronizability of $n \mapsto \seq{A026232}(n)$]
\label{rem:automata}
Proposition~\ref{prop:a189640-rec} shows that reading the base-$3$ expansion of $t \in \N$ least-significant-digit first, both $\seq{A189640}(t + 1)$ and the branching indicator $[\Pi_{66}(t + 2) \ne 2t + 2]$ are computed by the same $3$-state LSD-base-$3$ deterministic finite automaton with output (DFAO) having states $\{q_{\mathrm{wait}}, q_0, q_1\}$, start state $q_{\mathrm{wait}}$, transitions
\[
\delta(q_{\mathrm{wait}}, 0) = q_{\mathrm{wait}}, \quad \delta(q_{\mathrm{wait}}, 1) = q_0, \quad \delta(q_{\mathrm{wait}}, 2) = q_1, \quad \delta(q_0, d) = q_0, \quad \delta(q_1, d) = q_1,
\]
and output $\tau(q_1) = 1$, $\tau(q_{\mathrm{wait}}) = \tau(q_0) = 0$ (i.e., $a(t + 1) = 1$ iff the least-significant nonzero ternary digit of $t$ is $2$).

By contrast, the \emph{index-to-position map} $n \mapsto \seq{A026232}(n) = \OnePos(\seq{A189640})(n)$ is \textbf{not} $3$-synchronized~\cite{allouche2003}. Indeed, let $c_1(m)$ denote the number of occurrences of the digit~$1$ in the base-$3$ representation of $m \in \N$. Induction on $m$ using Proposition~\ref{prop:a189640-rec} shows that the prefix sum $S(N) := \sum_{i=1}^N \seq{A189640}(i)$ satisfies the exact identity
\[
2\,S(N) = (N - 1) - c_1(N - 1) \qquad (\forall\, N \ge 1),
\]
because incrementing $m = N - 1$ by~$1$ increases $c_1$ by~$+1$ when the least-significant nonzero ternary digit of $m + 1$ is~$1$ ($a(m + 2) = 0$) and decreases $c_1$ by~$-1$ when that digit is~$2$ ($a(m + 2) = 1$). Now take $j \ge 1$ and put $N_{2j} := (3^{2j+1} + 3)/2$, so that $N_{2j} - 1 = (1^{2j}2)_3$ has least-significant digit~$2$ (hence $a(N_{2j}) = 1$) and $c_1(N_{2j} - 1) = 2j$. Then $N_{2j} = (1^{2j-1}20)_3$ is the $n_{2j}$-th $1$-position, where
\[
n_{2j} = S(N_{2j}) = \frac{3^{2j+1} + 1}{4} - j = \bigl(0\,\underbrace{20\,20\cdots 20}_{j\text{ blocks}}\,1\bigr)_3 - j.
\]
Suppose toward a contradiction that a finite automaton with $s$ states recognized the base-$3$ graph $\{(n)_3 \otimes (N)_3 : N = \OnePos(\seq{A189640})(n)\}$. Choose $j$ large enough that $j < 3^r$ while $j - r > s + 2$; then subtracting $j$ cannot borrow past the low $r + 1$ digits, so the middle $2(s+1)$ digits of $(n_{2j})_3 \otimes (N_{2j})_3$ form the constant pair-word $\bigl((2,1)(0,1)\bigr)^{s+1}$. By the pigeonhole principle on the $s + 1$ pair boundaries, the automaton repeats a state after $p \ge 1$ pairs ($2p$ digits). Pumping that loop once yields an accepted pair whose second track equals $(1^{2(j+p)-1}20)_3 = N_{2(j+p)}$, while its first track evaluates to $(3^{2(j+p)+1} + 1)/4 - j = n_{2(j+p)} + p \ne n_{2(j+p)}$, contradicting functionality of $n \mapsto \OnePos(\seq{A189640})(n)$.

Even in automata-theoretic decision procedures such as Walnut, one therefore cannot represent $\seq{A026232}(n)$ directly as a $3$-synchronized function of~$n$; one must pass through the $3$-automatic characteristic set $K \in \seq{A026231}$. Our proof bypasses automata libraries entirely: Proposition~\ref{prop:rec-match} matches the two ternary recurrences directly on~$\N$ via strong induction, and Lemma~\ref{lem:strict-mono-eq} lifts the resulting set equivalence directly to sequence equality.
\end{remark}

\section{Lean 4 Formalization and Verification Record}
\label{sec:lean}

The complete formalization resides in the repository \texttt{derekblair/math-proofs} under Lean~\texttt{v4.34.1} (\nolinkurl{leanprover/lean4:v4.34.1}) and Mathlib revision
{\small\nolinkurl{d13f23b723b8a846827a245b89c10fc7d3f11612}}:
\begin{itemize}
\item \nolinkurl{MathProofs/SeamB/C018/Statement.lean} (86 lines; SHA-256:\\
{\small\nolinkurl{cf8e092edee68db4fff56176b78927b04cfc70ba60e77d6f8b4af77ba6e0602d}}):\\
verbatim transcription of the eleven definitions and two theorem signatures from the frozen pin manifest \texttt{PIN-C018.md} v1 (\texttt{delta\_vs\_source: NONE}, SHA-256:\\
{\small\nolinkurl{da0d52a6d4966c8908f92a046556ccf952b73d25139d9c66b5a81e5d18ba3a69}}), placed in \texttt{namespace MathProofs.SeamB.C018}.
\item \nolinkurl{MathProofs/SeamB/C018/Solution.lean} (836 lines; SHA-256:\\
{\small\nolinkurl{e329768e9b6005d729ae07251ffd73d2f45744322bb4b10cef8b351fd92aba6b}}):\\
self-contained zero-\texttt{sorry} proof importing only \nolinkurl{Mathlib.Data.Nat.Nth}, repeating the pinned definitions block and theorem signatures with $0$ character differences against \texttt{Statement.lean}.
\item \nolinkurl{Mathproofs.lean} (91 lines; SHA-256:\\
{\small\nolinkurl{9ce3e5fe146b6957d69cbc23fbe9f02fa8b5f2462d923d9ab455feafa205f351}}):\\
single import line \texttt{import MathProofs.SeamB.C018.Solution} appended to the library root.
\end{itemize}

\begin{table}[ht]
\centering
\caption{Declaration map in \texttt{namespace MathProofs.SeamB.C018} (\nolinkurl{MathProofs/SeamB/C018/Solution.lean}).}
\label{tab:lean-map}
\smallskip
{\small
\begin{tabular}{@{}l l l@{}}
\toprule
Paper Result & Lean~4 Declaration & Status \\
\midrule
Morphism & \texttt{a189640Morph}, \texttt{a189640Iter}, \texttt{a189640} & Pinned \texttt{def} \\
Position enumerations & \texttt{a189640ZeroPos}, \texttt{a189640OnePos} & Pinned \texttt{def} \\
Left-right filling & \texttt{lrState}, \texttt{lrFill}, \texttt{a026166}, \texttt{a026142} & Pinned \texttt{def} \\
Derived targets & \texttt{a026231}, \texttt{a026232} & Pinned \texttt{def} \\
Proposition~\ref{prop:a189640-rec} & \texttt{a189640\_one}, \texttt{a189640\_three\_mul\_add\_\{one,two,three\}} & Exported, proved \\
Proposition~\ref{prop:a026142} & \texttt{a026142\_three\_mul\_\{add\_two,add\_one,cases\}} & Exported, proved \\
Theorem~\ref{thm:a026166} (residues) & \texttt{a026166\_three\_mul}, \texttt{\_add\_one}, \texttt{\_add\_two} & Exported, proved \\
Proposition~\ref{prop:rec-match} & \texttt{a026166\_ne\_two\_mul\_iff\_a189640} & Exported, proved \\
Proposition~\ref{prop:set-equiv} & \texttt{a026231\_pred\_iff} & Exported, proved \\
Theorem~\ref{thm:main}(a) & \texttt{a026231\_eq\_three\_mul\_onePos\_sub\_one} & Exported, proved \\
Theorem~\ref{thm:main}(b) & \texttt{a026231\_mod\_three} (pinned, guard $1 \le n$) & Exported, proved \\
Theorem~\ref{thm:main}(c) & \texttt{a189640\_onePos\_eq\_a026232} (pinned, guard $1 \le n$) & Exported, proved \\
Offset-$1$ calibration & \texttt{c018\_offset\_one\_calibration} & Exported, proved \\
\bottomrule
\end{tabular}
}
\end{table}

\noindent\textbf{Audit summary.}
\begin{enumerate}
\item \emph{Clean-slate module build:} \nolinkurl{lake build MathProofs.SeamB.C018.Solution} succeeds with exit code~$0$, $828$ jobs, $3.1\text{ s}$ build time ($4.30\text{ s}$ wall), and $0$ warnings.
\item \emph{Statement fidelity:} Automated diff confirms $0$ differing characters between \texttt{Statement.lean} and \texttt{Solution.lean} across all eleven pinned definitions and both pinned theorem signatures.
\item \emph{Lexical trust-escape scan:} $0$ raw hits and $0$ comment/string-stripped hits for \texttt{sorry}, \texttt{admit}, \texttt{sorryAx}, \texttt{axiom}, \texttt{native\_decide}, \texttt{ofReduceBool}, \texttt{ofReduceNat}, \texttt{implemented\_by}, or \texttt{unsafe} in \texttt{Solution.lean}.
\item \emph{Axiom audit:} Both \texttt{\#print axioms} and a \texttt{CoreM} \texttt{Lean.collectAxioms} transitive check confirm that \texttt{a026231\_mod\_three}, \texttt{a189640\_onePos\_eq\_a026232}, and every exported theorem in Table~\ref{tab:lean-map} depend strictly on the standard Lean allowlist \texttt{[propext, Classical.choice, Quot.sound]} (\texttt{disallowed = []}).
\item \emph{Non-vacuity calibration:} \texttt{c018\_offset\_one\_calibration} verifies directly inside Lean's kernel that $\OnePos(\seq{A189640})(1) = 3$, $\seq{A026231}(1) = 8$, $\seq{A026232}(1) = 3$, $\Pi_{66}(8) = 6$, and $\Pi_{42}(8) = 5$, confirming $1$-based offset alignment without circularity.
\item \emph{Independent kernel replay:}
\texttt{lake env leanchecker --fresh --verbose} on module \nolinkurl{MathProofs.SeamB.C018.Solution} completes with exit code~$0$.
\end{enumerate}

\noindent\textbf{Scope boundaries.}
This work resolves \texttt{C018} (\seq{A189640} $1$-positions $= \seq{A026232}$) and characterizes \seq{A026166} completely. It does not re-formalize the companion $0$-positions identity \texttt{C017} ($\ZeroPos(\seq{A189640}) = \seq{A026138}$, already proved by Dekking~\cite{dekking2020}), nor does it formalize R.~J.\ Mathar's June~2025 empirical comment on \seq{A026232} relating \seq{A026232} to a left shift of \seq{A026227}.

\begin{thebibliography}{99}
\raggedright

\bibitem{oeis189640}
C.\ Kimberling, \emph{OEIS sequence \seq{A189640}}: Fixed point of the morphism $0 \to 001$, $1 \to 101$. Submitted April~24, 2011; includes proof of shift relation to \seq{A080846} by F.~M.\ Dekking (October~26, 2019); live revision \#45 (January~22, 2026). \url{https://oeis.org/A189640}.

\bibitem{oeis026232}
C.\ Kimberling, \emph{OEIS sequence \seq{A026232}}: $a(n) = (1/3)(s(n) + 1)$, where $s = \seq{A026231}$. Live revision \#6 (June~24, 2025). \url{https://oeis.org/A026232}.

\bibitem{oeis026231}
C.\ Kimberling, \emph{OEIS sequence \seq{A026231}}: Numbers $k$ such that $\seq{A026166}(k) = \seq{A026142}(k) + 1$. Title corrected by S.~A.\ Irvine (September~22, 2019); live revision \#14 (July~4, 2021). \url{https://oeis.org/A026231}.

\bibitem{oeis026166}
C.\ Kimberling, \emph{OEIS sequence \seq{A026166}}: Left-right filling permutation with $h = \lfloor (n-1)/2 \rfloor$. Live revision \#35 (July~3, 2020). \url{https://oeis.org/A026166}.

\bibitem{oeis026142}
C.\ Kimberling (renamed by P.\ Munn, January~7, 2022), \emph{OEIS sequence \seq{A026142}}: Lexicographically earliest permutation of the positive integers such that for $n \ge 2$, $|a(n) - n| = \lfloor (a(n)+1)/2 \rfloor$ (left-right filling permutation with $h = \lfloor (n+1)/2 \rfloor$). Live revision \#16 (January~14, 2022). \url{https://oeis.org/A026142}.

\bibitem{oeis026138}
C.\ Kimberling, \emph{OEIS sequence \seq{A026138}}: Numbers $k$ such that $s(j) < s(k)$ for all $j < k$, where $s = \seq{A026136}$. Live revision \#17 (August~28, 2021). \url{https://oeis.org/A026138}.

\bibitem{oeis080846}
N.~J.~A.\ Sloane, \emph{OEIS sequence \seq{A080846}}: Fixed point of the morphism $0 \to 010$, $1 \to 011$, starting from $a(1) = 0$. Live revision \#81 (September~3, 2026). \url{https://oeis.org/A080846}.

\bibitem{dekking2020}
F.~M.\ Dekking, \emph{Permutations of $\N$ generated by left-right filling algorithms}, J.\ Integer Seq.\ \textbf{23} (2020), Article~20.4.7; preprint \texttt{arXiv:2001.08915v1} [math.CO] (January~24, 2020). \url{https://cs.uwaterloo.ca/journals/JIS/VOL23/Dekking/dekk5.html}.

\bibitem{allouche2003}
J.-P.\ Allouche and J.\ Shallit, \emph{Automatic Sequences: Theory, Applications, Generalizations}, Cambridge University Press, Cambridge, 2003.

\end{thebibliography}

\end{document}
